\documentclass{article}
\usepackage{amsmath,amssymb,amsthm}
\usepackage{algorithm}
\usepackage{algorithmic}
\usepackage{geometry}
\geometry{margin=1in}
\newtheorem{theorem}{Theorem}
\newtheorem{lemma}{Lemma}
\newcommand{\E}{\mathbb{E}}
\newcommand{\R}{\mathbb{R}}
\begin{document}

\section*{Heavy‑Ball Gradient Method (HBGM)}

\section*{Mathematical Formulation}
Let $f:\R^{d}\rightarrow\R$ be differentiable.  
Given an initial pair of points $x_{0},x_{-1}\in\R^{d}$, step size $\alpha>0$ and momentum parameter $\beta\in[0,1)$, the Heavy‑Ball updates are
\[
\boxed{
x_{k+1}=x_{k}+ \beta\bigl(x_{k}-x_{k-1}\bigr)-\alpha\nabla f(x_{k}),\qquad k=0,1,2,\dots
}
\tag{HB}
\]
The term $\beta(x_{k}-x_{k-1})$ is the classical momentum (``velocity'') that uses the previous iterate difference, not a look‑ahead gradient.

\section*{Pseudocode}
\begin{algorithm}[H]
\caption{Heavy‑Ball Gradient Method}
\begin{algorithmic}[1]
\REQUIRE Initial point $x_{0}\in\R^{d}$, previous point $x_{-1}=x_{0}$, step size $\alpha>0$, momentum $\beta\in[0,1)$, max iterations $K$.
\FOR{$k=0$ \TO $K-1$}
    \STATE Compute gradient $g_{k}= \nabla f(x_{k})$.
    \STATE Update velocity $v_{k}= \beta (x_{k}-x_{k-1})$.
    \STATE $x_{k+1}= x_{k}+v_{k}-\alpha g_{k}$.
\ENDFOR
\RETURN $x_{K}$
\end{algorithmic}
\end{algorithm}

\section*{Dry Run Example}
Consider the univariate quadratic $f(w)=(w-3)^{2}$.  
Then $\nabla f(w)=2(w-3)$.  Choose $\alpha=0.1$, $\beta=0.5$, $x_{0}=0$, and set $x_{-1}=x_{0}=0$.

\begin{center}
\begin{tabular}{c|c|c|c|c}
$k$ & $x_{k}$ & $g_{k}= \nabla f(x_{k})$ & $v_{k}= \beta(x_{k}-x_{k-1})$ & $x_{k+1}=x_{k}+v_{k}-\alpha g_{k}$\\\hline
0 & $0$ & $2(0-3)=-6$ & $0.5(0-0)=0$ & $0+0-0.1(-6)=0.6$\\
1 & $0.6$ & $2(0.6-3)=-4.8$ & $0.5(0.6-0)=0.3$ & $0.6+0.3-0.1(-4.8)=1.38$\\
2 & $1.38$ & $2(1.38-3)=-3.24$ & $0.5(1.38-0.6)=0.39$ & $1.38+0.39-0.1(-3.24)=2.104$\\
\end{tabular}
\end{center}
Objective values: $f(0)=9$, $f(0.6)=5.76$, $f(1.38)=2.6244$, $f(2.104)=0.803$ – the method reduces the function value quickly.

\newpage
\section*{Assumptions}
\begin{enumerate}
\item \textbf{Smoothness.} $f$ is $L$‑smooth, i.e.
\[
\|\nabla f(x)-\nabla f(y)\|\le L\|x-y\|,\qquad\forall x,y\in\R^{d}.
\tag{A1}
\]
Equivalently,
\[
f(y)\le f(x)+\langle\nabla f(x),y-x\rangle+\frac{L}{2}\|y-x\|^{2}.
\tag{A1'}
\]

\item \textbf{Lower boundedness.} $f$ has a finite infimum $f^{\inf}>-\infty$.

\item \textbf{Parameter choice.} The stepsize $\alpha$ and momentum $\beta$ satisfy
\[
0<\alpha\le \frac{1}{L},\qquad 0\le\beta\le\sqrt{1-\alpha L}.
\tag{A2}
\]
The upper bound on $\beta$ guarantees the Lyapunov function used in the proof is non‑increasing.

\item \textbf{Initial velocity.} $x_{-1}=x_{0}$ (zero initial momentum). This simplifies the analysis without loss of generality.
\end{enumerate}

\section*{Main Convergence Theorem}
\begin{theorem}[Non‑convex convergence of Heavy‑Ball]\label{thm:hb}
Under assumptions (A1)–(A2), let $\{x_{k}\}$ be generated by \eqref{HB}. Define the averaged gradient norm
\[
\bar{g}_{K}:=\frac{1}{K}\sum_{k=0}^{K-1}\|\nabla f(x_{k})\|^{2}.
\]
Then for any $K\ge 1$,
\[
\bar{g}_{K}\;\le\;
\frac{2\bigl(f(x_{0})-f^{\inf}\bigr)}{\alpha K}
\;+\;
\frac{2\beta}{\alpha K}\|x_{0}-x_{-1}\|^{2}.
\tag{T1}
\]
In particular, when $x_{-1}=x_{0}$ the second term disappears and we obtain the rate
\[
\boxed{\displaystyle
\frac{1}{K}\sum_{k=0}^{K-1}\|\nabla f(x_{k})\|^{2}=O\!\left(\frac{1}{K}\right)}.
\]
\end{theorem}

\section*{Theoretical Properties}
\begin{itemize}
\item \textbf{Stability.} Condition \eqref{A2} is tight for the quadratic case; violating it may cause divergence (oscillations) because the Lyapunov function ceases to be monotone.

\item \textbf{Hyper‑parameter sensitivity.} Larger $\beta$ accelerates progress when the objective is close to a quadratic basin, but the bound $\beta\le\sqrt{1-\alpha L}$ caps $\beta$ for a given $\alpha$.

\item \textbf{Comparison to Gradient Descent (GD).} GD corresponds to $\beta=0$, yielding $\bar{g}_{K}\le 2\bigl(f(x_{0})-f^{\inf}\bigr)/(\alpha K)$. The extra momentum term does not worsen the $O(1/K)$ rate, yet empirically often reduces the constant factor.

\item \textbf{Special case – convex smooth $f$.} The same proof gives the classical $O(1/K)$ function‑value convergence for convex $f$ (see Polyak, 1964).
\end{itemize}

\section*{Discussion}
The theorem shows that the Heavy‑Ball method attains the same worst‑case gradient‑norm decay as plain GD for non‑convex $L$‑smooth functions, provided the stepsize and momentum respect \eqref{A2}. The proof relies on a simple Lyapunov (energy) function that couples function values and the kinetic term $\|x_{k}-x_{k-1}\|^{2}$. The resulting bound is deterministic; no stochasticity is required.

\newpage
\section*{Preliminaries (Definitions and Lemmas)}
\begin{lemma}[Descent Lemma]\label{lem:descent}
If $f$ is $L$‑smooth then for any $x,y\in\R^{d}$
\[
f(y)\le f(x)+\langle\nabla f(x),y-x\rangle+\frac{L}{2}\|y-x\|^{2}.
\tag{L1}
\]
\end{lemma}
\begin{proof}
Standard from the integral form of the gradient; omitted for brevity. \qedhere
\end{proof}

\begin{lemma}[Quadratic bound for the momentum term]\label{lem:beta}
For any $a,b\in\R^{d}$ and $\beta\in[0,1)$,
\[
\|a+\beta b\|^{2}\le (1+\beta)\|a\|^{2}+(1+\beta^{-1})\beta^{2}\|b\|^{2}.
\tag{L2}
\]
\end{lemma}
\begin{proof}
Expand $\|a+\beta b\|^{2}= \|a\|^{2}+2\beta\langle a,b\rangle+\beta^{2}\|b\|^{2}$ and apply Cauchy–Schwarz together with $2\beta|\langle a,b\rangle|\le \beta(\|a\|^{2}+\|b\|^{2})$. Rearranging yields (L2). \qedhere
\end{proof}

\section*{Main Proof (Step‑by‑Step Derivation)}
We prove Theorem~\ref{thm:hb}. Throughout the proof we denote $g_{k}:=\nabla f(x_{k})$.

\bigskip
\noindent\textbf{Step 1.} Define the Lyapunov function
\[
\Phi_{k}:=f(x_{k})+\frac{\beta}{2\alpha}\|x_{k}-x_{k-1}\|^{2},\qquad k\ge0.
\tag{S1}
\]

\noindent\textbf{Step 2.} Using the update \eqref{HB} write $x_{k+1}-x_{k}= \beta (x_{k}-x_{k-1})-\alpha g_{k}$ and apply Lemma~\ref{lem:descent} with $x=x_{k}$, $y=x_{k+1}$:
\[
f(x_{k+1})\le f(x_{k})+\langle g_{k},x_{k+1}-x_{k}\rangle
+\frac{L}{2}\|x_{k+1}-x_{k}\|^{2}.
\tag{S2}
\]

\noindent\textbf{Step 3.} Substitute $x_{k+1}-x_{k}= \beta (x_{k}-x_{k-1})-\alpha g_{k}$ into (S2):
\[
\begin{aligned}
f(x_{k+1})
&\le f(x_{k})+\langle g_{k},\beta (x_{k}-x_{k-1})-\alpha g_{k}\rangle
+\frac{L}{2}\|\beta (x_{k}-x_{k-1})-\alpha g_{k}\|^{2}\\
&= f(x_{k})-\alpha\|g_{k}\|^{2}
+\beta\langle g_{k},x_{k}-x_{k-1}\rangle
+\frac{L}{2}\bigl(\beta^{2}\|x_{k}-x_{k-1}\|^{2}
-2\alpha\beta\langle g_{k},x_{k}-x_{k-1}\rangle
+\alpha^{2}\|g_{k}\|^{2}\bigr).
\end{aligned}
\tag{S3}
\]

\noindent\textbf{Step 4.} Collect the terms involving $\langle g_{k},x_{k}-x_{k-1}\rangle$:
\[
\beta\langle g_{k},x_{k}-x_{k-1}\rangle
-\alpha L\beta\langle g_{k},x_{k}-x_{k-1}\rangle
= \beta\bigl(1-\alpha L\bigr)\langle g_{k},x_{k}-x_{k-1}\rangle.
\tag{S4}
\]

\noindent\textbf{Step 5.} Use the elementary inequality $2ab\le \eta a^{2}+b^{2}/\eta$ with $\eta>0$ to bound the mixed term in (S4). Choose $\eta=\alpha L$:
\[
\beta\bigl(1-\alpha L\bigr)\langle g_{k},x_{k}-x_{k-1}\rangle
\le \frac{\beta(1-\alpha L)}{2}\Bigl(\alpha L\|x_{k}-x_{k-1}\|^{2}
+\frac{1}{\alpha L}\|g_{k}\|^{2}\Bigr).
\tag{S5}
\]

\noindent\textbf{Step 6.} Plug (S5) into (S3) and simplify the coefficients of $\|g_{k}\|^{2}$ and $\|x_{k}-x_{k-1}\|^{2}$.
After algebraic manipulation (use $\alpha\le 1/L$ and $0\le\beta\le\sqrt{1-\alpha L}$) we obtain
\[
f(x_{k+1})\le f(x_{k})-\frac{\alpha}{2}\|g_{k}\|^{2}
+\frac{\beta}{2\alpha}\bigl(\beta^{2}- (1-\alpha L)\bigr)\|x_{k}-x_{k-1}\|^{2}.
\tag{S6}
\]

\noindent\textbf{Step 7.} Add $\dfrac{\beta}{2\alpha}\|x_{k+1}-x_{k}\|^{2}$ to both sides of (S6). Recall $x_{k+1}-x_{k}= \beta (x_{k}-x_{k-1})-\alpha g_{k}$, expand the square and use Lemma~\ref{lem:beta} with $a=-\alpha g_{k}$ and $b=\beta (x_{k}-x_{k-1})$:
\[
\|x_{k+1}-x_{k}\|^{2}\le (1+\beta)\alpha^{2}\|g_{k}\|^{2}
+(1+\beta^{-1})\beta^{2}\|x_{k}-x_{k-1}\|^{2}.
\tag{S7}
\]

\noindent\textbf{Step 8.} Multiply (S7) by $\dfrac{\beta}{2\alpha}$ and add to (S6):
\[
\Phi_{k+1}\le \Phi_{k}-\underbrace{\Bigl(\frac{\alpha}{2}-\frac{\beta(1+\beta)}{2}\alpha\Bigr)}_{\displaystyle\ge 0}
\|g_{k}\|^{2}
+\underbrace{\frac{\beta}{2\alpha}\Bigl[\beta^{2}-(1-\alpha L)+(1+\beta^{-1})\beta^{2}\Bigr]}_{\displaystyle\le 0}
\|x_{k}-x_{k-1}\|^{2}.
\tag{S8}
\]

\noindent\textbf{Step 9.} By the parameter restriction (A2) we have $ \beta^{2}\le 1-\alpha L$, which makes the second bracket non‑positive. Hence the kinetic term can be dropped, yielding the monotone inequality
\[
\Phi_{k+1}\le \Phi_{k}-\frac{\alpha}{2}\|g_{k}\|^{2}.
\tag{S9}
\]

\noindent\textbf{Step 10.} Telescoping (S9) from $k=0$ to $K-1$ gives
\[
\Phi_{K}\le \Phi_{0}-\frac{\alpha}{2}\sum_{k=0}^{K-1}\|g_{k}\|^{2}.
\tag{S10}
\]

\noindent\textbf{Step 11.} Since $f$ is lower bounded,
$\Phi_{K}=f(x_{K})+\dfrac{\beta}{2\alpha}\|x_{K}-x_{K-1}\|^{2}\ge f^{\inf}$.
Thus,
\[
\frac{\alpha}{2}\sum_{k=0}^{K-1}\|g_{k}\|^{2}
\le \Phi_{0}-f^{\inf}
= f(x_{0})-f^{\inf}+\frac{\beta}{2\alpha}\|x_{0}-x_{-1}\|^{2}.
\tag{S11}
\]

\noindent\textbf{Step 12.} Divide both sides of (S11) by $\alpha K$ and rearrange:
\[
\frac{1}{K}\sum_{k=0}^{K-1}\|g_{k}\|^{2}
\le \frac{2\bigl(f(x_{0})-f^{\inf}\bigr)}{\alpha K}
+\frac{2\beta}{\alpha K}\|x_{0}-x_{-1}\|^{2}.
\tag{S12}
\]
This is exactly the bound \eqref{T1} of Theorem~\ref{thm:hb}.

\noindent\textbf{Step 13.} When $x_{-1}=x_{0}$ the second term vanishes, yielding the $O(1/K)$ rate.

\hfill $\square$

\section*{Rate Derivation (With Constants)}
From \eqref{T1} we can read the explicit constants:
\[
C_{1}= \frac{2\bigl(f(x_{0})-f^{\inf}\bigr)}{\alpha},\qquad
C_{2}= \frac{2\beta}{\alpha}\|x_{0}-x_{-1}\|^{2},
\]
so that
\[
\bar{g}_{K}\le \frac{C_{1}}{K}+\frac{C_{2}}{K}.
\]
If $x_{-1}=x_{0}$, $C_{2}=0$ and the bound simplifies to $\bar{g}_{K}\le C_{1}/K$.

\section*{Special Cases}
\begin{enumerate}
\item \textbf{Quadratic $f(x)=\tfrac12 x^{\top}Qx$ with $0\prec mI\preceq Q\preceq LI$.}  
The Heavy‑Ball recursion can be diagonalised; the condition $\beta\le\sqrt{1-\alpha L}$ is also necessary for spectral radius $<1$, confirming the tightness of (A2).

\item \textbf{Gradient Descent.} Setting $\beta=0$ in \eqref{HB} reduces the theorem to the classical GD bound $\frac{1}{K}\sum\|\nabla f(x_{k})\|^{2}\le 2(f(x_{0})-f^{\inf})/(\alpha K)$.

\item \textbf{Large momentum.} If $\beta>\sqrt{1-\alpha L}$, the coefficient in (S9) may become negative, and the Lyapunov function may increase, leading to possible divergence – a phenomenon observed empirically.
\end{enumerate}

\end{document}